\documentclass[12pt]{article}
\usepackage[utf8]{inputenc}
\usepackage[T1]{fontenc}
\usepackage{hyperref}
\title{Martuyhunira Language Resources for K-12 Curriculum Development}
\author{Compiled by Claude Sonnet 5 (claude-sonnet-5) for Andr\'as Kornai}
\date{}
\begin{document}
\maketitle

\section{Introduction}
Martuyhunira (ISO 639-3 vma, Glottocode mart1255) is the Ethnologue spelling. The usual spelling is Martuthunira. It is a Ngayarta language of the Pama-Nyungan family, originally spoken in the south-western Pilbara of Western Australia, around the Fortescue and Robe River area. Sources differ on the end of fluent speech: Wikipedia reports that the last fluent speaker, Algy Paterson, died on 6 August 1995, and Wangka Maya reports an estimate of fewer than five speakers in 2004. The language is no longer spoken fluently.

\section{Primary Educational Resources}
I found no school primers or textbooks for Martuthunira. Wangka Maya Pilbara Aboriginal Language Centre supports Pilbara languages generally.
\begin{itemize}
\item Wangka Maya, Martuthunira page: \url{https://www.wangkamaya.org.au/pilbara-languages/martuthunira}
\item Dixon and Deak, chapter on the Wangka Maya language centre as a revitalisation strategy (\textit{Re-Awakening Languages}, Sydney University Press): \url{https://open.sydneyuniversitypress.com.au/9781920899554/rl-language-centre-as-language-revitalisation-strategy-a-case-study-from-the-pilbara.html}
\end{itemize}

\section{Linguistic and Reference Materials}
\begin{itemize}
\item Alan Dench, \textit{Martuthunira: a language of the Pilbara region of Western Australia} (reference grammar, Pacific Linguistics, ANU repository): \url{https://openresearch-repository.anu.edu.au/handle/1885/10733}
\item Martuthunira language, Wikipedia: \url{https://en.wikipedia.org/wiki/Martuthunira_language}
\item AIATSIS Austlang entry W35: \url{https://collection.aiatsis.gov.au/austlang/language/w35}
\item Language Description Heritage entry for vma: \url{https://ldh.clld.org/languages/vma}
\end{itemize}

\section{Religious and Community Resources}
\begin{itemize}
\item Mardudunera people, Wikipedia: \url{https://en.wikipedia.org/wiki/Mardudunera_people}
\end{itemize}

\section{Digital and Audio Resources}
I found no audio or video resources in these searches. The AIATSIS catalogue entry above is a starting point for archived recordings.

\section{Curriculum Development Context}
Dench's grammar is based largely on the speech of one man, Algy Paterson, so documentation exists but is narrow. A heritage and revitalisation curriculum would have to be built with the Pilbara community and Wangka Maya, using the grammar and archived material as source. I found no existing school programme for Martuthunira.

\end{document}